\section{Results for $C_5$}\label{sec:C5}

The results for $C_5$ when the graph has constant arboricity are very similar to those for $C_4$ (when the arboricity $\alpha$ is a constant), and hence we only provide a sketch in this section.

We start with the upper bound (as stated in Theorem~\ref{thm:C4-5} in the introduction).
\begin{theorem}\label{thm:C5-ub}
There exists a one-sided error algorithm for testing $C_5$-freeness in graphs with  arboricity at most $\alpha = O(1)$ over $n$ vertices whose query complexity is $\widetilde{O}(n^{1/4}/\eps^{4})$. 
% \commentt{I added another factor of $1/\eps$ because of the extra depth of the BFS/random walks}
\end{theorem}

The algorithm is an adaptation of the algorithm for testing $C_4$-freeness.
 While the algorithm for testing $C_4$-freeness attempted to find a $C_4$ by obtaining two length-$2$ edge-disjoint paths from a vertex $v$ to a vertex $v'$, the algorithm described next
 %Algorithm~\ref{alg:C5} 
 attempts to find a $C_5$ by obtaining two edge-disjoint paths from a vertex $v$ to a vertex $v'$ such that one of these paths is of length $2$ and the other of length $3$.

\medskip\noindent
\textsf{Test-$C_5$-freeness}($n,\epsilon,\alpha$)\label{alg:C5}
\begin{enumerate}
\item Let $\theta_0 = 4\alpha/\epsilon$ and $\theta_1 =  c_1\sqrt{n}/\eps$ (for a sufficiently large constant $c_1$).
\label{step:set_params_C5}
\item Repeat the following %$t=100/\epsilon$ 
$t = \Theta(1/\eps)$ times:
  \begin{enumerate}
     \item Select an edge $e$ by calling the procedure \textsf{Select-an-edge}($\alpha,\epsilon$) (as provided in Section~\ref{subsec:C4-ub}).  If it
     % fails to return an edge, then return `Accept' (and terminate).
     does not return an edge, then continue to the next iteration.
     \item Select an endpoint $v$ of $e$ by flipping a fair coin.
     \item \label{step:random-neighbors-C5} If $d(v) \leq \theta_1$, then select %$s=64\sqrt{d(v)/\eps}$ 
     $s_1= \Theta(\sqrt{d(v)/\eps})$ 
     random neighbors of $v$, and for each neighbor $u$ such that $d(u) \leq \theta_0$, query all the neighbors of $u$, and all the neighbors $u'$ of $u$ for which $d(u) \leq \theta_0$.
     That is, from each of the $s_1$ selected neighbors of $v$, perform a restricted BFS of depth 2, where only vertices with degree at most $\theta_0$ are further explored.
     \item \label{step:random-walk-C5} Otherwise ($d(v) > \theta_1$), perform 
     $s_2 = \Theta(\eps^{-3}\cdot\sqrt{n\log n/\theta_1})$ ~($=\tilde{\Theta}(n^{1/4}/\eps^3)$)
    %  $s=(10^5/\eps^3)\cdot \sqrt{2\log n (2m/\theta_1)}$
     random walks of length $3$ starting from $v$. %\commentt{Or should we select $s$ random neighbors of $v$ and then perform a BFS of depth-$2$ from each as in the first case? If so, we may be able to unify the two items (though then it will be less similar). }
     \item If a $C_5$ is detected, then return it, `Reject' and terminate.
  \end{enumerate}
  \item Return `Accept'.
\end{enumerate}
Since the algorithm rejects a graph only if it observes a $C_5$, it always accepts $C_5$-free graphs. Hence, consider a graph $G$  that is $\eps$-far from being $C_5$-free. 
Similarly to the analysis of the algorithm for testing $C_4$-freeness, we let $\calC$ be a maximal set of edge-disjoint $C_5$s such that no $C_5$ in this set contains an edge between two vertices with degree greater than $\theta_0$. We then define $\calC_1$ and $\calC_2$ to be the subsets of $\calC$ containing $C_5$s with no edge between two vertices with degree greater than $\theta_0$ and at most one (respectively, exactly two) vertices with degree greater then $\theta_1$. Here too, either $|\calC_1|\geq \eps m/16$ or $|\calC_2|\geq \eps m/16$.

If $|\calC_1|\geq \eps m/16$, then by applying a coloring process as in the proof of Theorem~\ref{thm:C4-ub}, we obtain a set $\calC'_1\subseteq \calC_1$
in which every vertex $v$ participating in some $C_5$ in $\calC'_1$ actually participates in  $\Omega(\eps d(v))$ such $C_5$s. Since every $5$-cycle $\rho\in \calC'_1$ contains at most one vertex of degree greater than $\theta_1$, we can label the vertices in $\rho$ by $v_i(\rho)$, $i=0,1,\dots,4$ where $v_2(\rho)$ is the highest degree vertex, such that the following holds. First, $v_0(\rho) \leq \theta_1$, and second, $v_1(\rho)$, $v_3(\rho)$ and $v_4(\rho)$ all have degree at most $\theta_0$. This implies that if the algorithm selects a vertex $v \in V_0(\calC_0') = \bigcup_{\rho\in \calC'_1} \{v_0(\rho)\}$ (which occurs with high constant probability), then with high constant probability the algorithm will detect a $C_5$ 
% in $\calC_1'$ that $v$ participates in 
in Step~\ref{step:random-neighbors-C5}.

If $|\calC_2| \geq \eps m/4$, then we apply a similar coloring process to obtain a set $\calC'_1\subseteq \calC_2$
in which every vertex $v$ participating in some $C_5$ in $\calC'_2$ actually participates in  $\Omega(\eps d(v))$ such $C_5$s. Since every $5$-cycle $\rho\in \calC'_2$ contains two vertices, $v$ and $v'$ of degree greater than $\theta_1$, there is one length-2 path between $v$ and $v'$ that passes through a vertex with degree at most $\theta_0$, and one length-3 path that passes through two vertices with degree at most $\theta_0$. An analysis similar to the given in the proof of Claim~\ref{clm:C4-ub2} shows that in this case (that $|\calC_2| \geq \eps m/4$), with high constant probability, a $C_5$ will be detected in Step~\ref{step:random-walk-C5} of the algorithm.
\ForFuture{T; elaborate a bit more}


\medskip\noindent
We now turn to the lower bound.
\begin{theorem} \label{thm:C5-const-alpha-lb}
Testing $C_5$-freeness in constant-arboricity graphs over $n$ vertices requires $\Omega(n^{1/4})$ queries for constant $\epsilon$.
\end{theorem}
%\commentt{Need to double check that indeed works for two-sided. If problematic, then state for one-sided}

The high-level idea of the proof of Theorem~\ref{thm:C5-const-alpha-lb} is very similar to that of Theorem~\ref{thm:C4-const-alpha-lb}. The main difference is that in both $G_0$ and $G_1$ (based on which the distributions $\calD_0$ and $\calD_1$ are defined), 
some of the vertices in $X_b$ (for $b\in \{0,1\}$) are connected by an edge (so that the graphs are not bipartite),
and for each such edge $(v_1,v_2)$, there are two vertices $u_1$ and $u_2$ in $Y_b$, such that $u_1$ neighbors $v_1$ and $u_2$ neighbors $v_2$.

In $G_0$, every pair of vertices $u_1,u_2$ in $Y_0$ either neighbor a unique vertex $v_0$ in $X_0$ or there is a unique edge $(v_1,v_2)$ between a pair of vertices in $X_0$ such that $u_1$ neighbors $v_1$ and $u_2$ neighbors $v_2$.
On the other hand, in $G_1$, edge-disjoint $C_5$s are created by having the same pair $u_1,u_2$ in $Y_1$ neighbor both a single vertex $v_0$ and the endpoints of an edge $(v_1,v_2)$ as described above, so that $(v_0,u_1,v_1,v_2,u_2)$ is a $C_5$.


